\section{Scalar curvature and electromagnetic comparison}\label{ch:cmp:section}

Fix the prepared datum, then $\epsilon$ and $n$, and take the smooth limiting
solution supplied by Proposition~\ref{ch:ell:neutral-input}.
Throughout this section, derivatives, tensor contractions and cross products
without a metric subscript are taken with respect to $g$.  Put
\begin{equation}\label{ch:cmp:metrics}
 \bar g=g+l^2\,df\otimes df,\qquad
 \hat g=e^{4t}\bar g,\qquad
 \widetilde g=e^{4\epsilon}\hat g.
\end{equation}
The scalar and boundary identities of the neutral construction apply
without changing the total constraint densities. We first identify their
notation here, then use the physical matter margin to retain the total
charge in the comparison metric.

\subsection{The scalar identity}

\begin{lemma}[Scalar identity]\label{ch:cmp:scalar}
With the definitions of Section~\ref{ch:ell:section}, including the
quadratic form in Equation~\eqref{ch:ell:quadratic}, one has
\begin{equation}\label{ch:cmp:scalar-identity}
 \frac12 e^{4t}R_{\hat g}
 =8\pi\bigl(\mu+J(w)\bigr)+\cT+w(F)-F\tr K
       -u^{-1}\Div V.
\end{equation}
Here $\mu,J$ are the total constraint densities of $(g,K)$.
The quadratic form $\cT$ extends smoothly through $\nabla f=0$ and is
nonnegative everywhere.
\end{lemma}

\begin{proof}
Apply Proposition~\ref{n3:deformation:identity}. The variables
$l,L,p,D,d,w,a,v,A_d,A_\chi,Z,u,H^f,F,V$ in
Equations~\eqref{ch:ell:profiles}--\eqref{ch:ell:variables} are those
of that proposition. Its tensor $S$ is our $S_f$, its axial component
$q$ is our $k$, and its transverse and axial derivatives of $t$ are
$y_t,x_t$. Its logarithmic derivative $p_L=L'/L$ is $p+2$.
In particular the graph and conformal metrics, momentum convention,
and full flux agree. The identity holds for arbitrary smooth $f,t$;
no elliptic equation or constraint beyond the definitions of $\mu,J$
is assumed in its proof.

Equation~\eqref{n3:deformation:square-completion} is exactly
Equation~\eqref{ch:ell:quadratic} under this correspondence. It proves
nonnegativity, since $p\ge0$ and $0<d\le1$. The invariant formula
\eqref{n3:deformation:quadratic-invariant} proves smoothness at
$\nabla f=0$, where its value is
\[
 \cT=\tfrac12\bigl(|S_f|^2+(\tr S_f)^2\bigr)
                          +4(p+1)|\nabla t|^2.
\]
Thus the shared identity gives \eqref{ch:cmp:scalar-identity}, with
the momentum term $+8\pi J(w)$.
\end{proof}

\begin{lemma}[Weighted coercivity]\label{ch:cmp:coercivity}
At every point, with any axial direction at zero slope,
\begin{equation}\label{ch:cmp:coercivity-estimate}
 |T|^2+|M|^2+dk^2+F^2+|y_t|^2+dx_t^2
                  \le13(1+n)^2\cT.
\end{equation}
The constant is independent of the solution, $n$, the truncation radius,
and the size of $|\nabla f|$.
\end{lemma}

\begin{proof}
Use the square completion
\eqref{n3:deformation:square-completion}, identified above with
\eqref{ch:ell:quadratic}. Its last coefficient is at least
$2d/[3(1+p)^2]$ and at least $2\chi^2/3$. Its transverse-gradient
coefficient is at least $3(p+1)$. Comparing these terms and the
shifted squares gives
\[
 \begin{gathered}
 dk^2\le\tfrac32(1+p)^2\cT,\qquad
 \chi^2k^2\le\tfrac32\cT,\qquad
 |y_t|^2\le\frac{\cT}{3(p+1)},\\
 F^2\le3\cT,\qquad
 |M|^2\le\tfrac83(p+1)\cT,\qquad |T|^2\le4\cT.
 \end{gathered}
\]
Here the shifted $F$ and $M$ estimates use
$|A+B|^2\le2|A|^2+2|B|^2$; for $T$ write
$\tr T=(F-\chi k/3)-2\chi k/3$. The axial-gradient square gives
\[
 dx_t^2\le\frac{\cT}{2(p+1)}+\frac{\chi^2vk^2}{8d}
 \le\left(\frac1{2(p+1)}+\frac16\right)\cT
 \le\tfrac23\cT,
\]
because the ratio of its second coefficient to the last square
coefficient is $3v/[2(9-d)]\le1/6$. Summing and using
$0\le p\le n$ proves \eqref{ch:cmp:coercivity-estimate}.
At zero slope the same estimates hold for every axial decomposition
of the direction-independent value of $\cT$.
\end{proof}

\subsection{The inner boundary}

\begin{lemma}[Boundary identity and sign]\label{ch:cmp:boundary}
On every face of $B$, with $\nu$ pointing into $X$ and
$P_B=\tr_BK$,
\begin{equation}\label{ch:cmp:boundary-identity}
 \frac{V_\nu}{u}
    =d(H+4\partial_\nu t)-H+w_\nu(F-P_B).
\end{equation}
For a solution of Equation~\eqref{ch:ell:system}, the corresponding mean
curvatures satisfy
\begin{equation}\label{ch:cmp:negative-boundary}
 H_{\hat g}=-e^{-2t}\sqrt d\,N_B<0,\qquad
 H_{\widetilde g}=e^{-2\epsilon}H_{\hat g}<0.
\end{equation}
Both mean curvatures use the normal toward the designated end.
\end{lemma}

\begin{proof}
The assigned height is constant on each face, and the normal into $X$
is the normal used in Lemma~\ref{n3:deformation:boundary-identity}.
Its flux, trace $F$, and tangential trace $P_B$ agree with ours by the
correspondence in Lemma~\ref{ch:cmp:scalar}. Applying
Lemma~\ref{n3:deformation:boundary-identity} gives
\eqref{ch:cmp:boundary-identity} and
\[
 H_{\hat g}=e^{-2t}\sqrt d\,(H+4\partial_\nu t).
\]
The boundary equation in \eqref{ch:ell:system} now implies
$d(H+4\partial_\nu t)=-N_Bd$. Since $d>0$, substitution gives
\eqref{ch:cmp:negative-boundary}; the last equality is the
mean-curvature scaling under $\widetilde g=e^{4\epsilon}\hat g$.
\end{proof}

\subsection{Retaining both total charges}

\begin{proposition}[Charged comparison]\label{ch:cmp:charged-comparison}
For every limiting solution furnished by Proposition~\ref{ch:ell:neutral-input}, there are
smooth closed two-forms $\widetilde\alpha_1,\widetilde\alpha_2$ on $X$
whose end fluxes are $4\pi Q_E,4\pi Q_B$, respectively, such that the
fields $\widetilde{\cE}=\widetilde X_1$, $\widetilde{\cB}=\widetilde X_2$ defined by
$\iota_{\widetilde X_i}dV_{\widetilde g}=\widetilde\alpha_i$ satisfy
\begin{equation}\label{ch:cmp:charged-scalar}
 R_{\widetilde g}\ge
       2\bigl(|\widetilde X_1|_{\widetilde g}^2
                         +|\widetilde X_2|_{\widetilde g}^2\bigr),
 \qquad \Div_{\widetilde g}\widetilde X_i=0.
\end{equation}
Moreover $\widetilde g\ge g$ as quadratic forms, $X$ is complete up to
$B$ for $\widetilde g$, and $H_{\widetilde g}<0$ on $B$.
\end{proposition}

\begin{proof}
Substituting Equation~\eqref{ch:ell:system} into
Equation~\eqref{ch:cmp:scalar-identity}, and using $w(h)=\tau_0a\sigma$, gives
\begin{align}
 \frac12e^{4t}R_{\hat g}
 ={}&8\pi(\mu+J(w))+(1-\delta_0)\cT
            +(1-\delta_0)\tau_0a\sigma\notag\\
    &+w(C)-F\tr K-\rho+m_0(t)\cP.\label{ch:cmp:scalar-endpoint}
\end{align}
The physical matter condition says
\[
 8\pi(\mu+J(w))-I_g(w)
 =8\pi(\mu_m+J_m(w))\ge M_0,
\]
since $|w|<1$.  The compact height bound in Proposition~\ref{ch:ell:neutral-input} and
Equation~\eqref{ch:prep:offset-smallness} imply
$|w(C)|+|F\tr K|+\rho\le M_0/2$; outside the supports of $K$ and $C$
only the already controlled term $\rho$ remains.
All other terms in Equation~\eqref{ch:cmp:scalar-endpoint} are nonnegative.
Consequently
\begin{equation}\label{ch:cmp:physical-scalar-bound}
 \frac12e^{4t}R_{\hat g}\ge I_g(w).
\end{equation}

For $Q>0$, take the constant rotation
\[
 O=\frac1Q\begin{pmatrix}Q_E&Q_B\\-Q_B&Q_E\end{pmatrix}\in SO(2),
 \qquad \binom{X_e}{X_o}=O\binom{\cE}{\cB};
\]
for $Q=0$, use $O=I$.  Rotate the flux forms by the same matrix.
Their charges become $(Q,0)$, while $I_g$ is unchanged: the sum of
squared norms is invariant and $X_e\times X_o=\cE\times\cB$.
Let $\alpha=\iota_{X_e}dV_g$.  Completing one vector square gives
\begin{align}
 I_g(w)&=|X_o+w\times X_e|^2+|X_e|^2-|w\times X_e|^2\notag\\
       &\ge|X_e|^2-|w\times X_e|^2
         =|\alpha|_{\bar g}^2.\label{ch:cmp:field-square}
\end{align}
The last equality is a two-form norm identity.
In a positively oriented $g$-orthonormal frame with first vector
parallel to $w$, the inverse graph metric has eigenvalues $(d,1,1)$.
Writing $X_e=(x_1,x_2,x_3)$ gives
\[
 \alpha=x_1 e^2\wedge e^3+x_2 e^3\wedge e^1+x_3 e^1\wedge e^2,
 \qquad
 |\alpha|_{\bar g}^2=x_1^2+d(x_2^2+x_3^2)
                       =|X_e|^2-|w\times X_e|^2.
\]
At $w=0$ this formula is independent of the frame.
This is the electric graph transport of \cite[Section~2 and Appendix~A]{DisconziKhuri2012},
expressed through its flux form.

Retain the rotated pair $(\alpha,0)$ and define the new pair by the inverse
rotation
\begin{equation}\label{ch:cmp:new-flux-forms}
 \binom{\widetilde\alpha_1}{\widetilde\alpha_2}
           =O^{\mathsf T}\binom{\alpha}{0}.
\end{equation}
This discards a form of zero \emph{total} charge.  Since $O$ is constant,
the new forms are closed, have separately the original total charges,
and satisfy $|\widetilde\alpha_1|_h^2+|\widetilde\alpha_2|_h^2=|\alpha|_h^2$
for every metric $h$.
No assertion about the individual component charges is needed.
Flux conservation across every enclosing cut follows from Stokes' Theorem.

Finally set $s_0=t+\epsilon\ge0$, using the floor in
Proposition~\ref{ch:ell:neutral-input}. Scalar curvature under constant scaling and the conformal norm
of a two-form give
\[
 R_{\widetilde g}=e^{-4\epsilon}R_{\hat g},\qquad
 |\alpha|_{\widetilde g}^2=e^{-8s_0}|\alpha|_{\bar g}^2.
\]
Equations~\eqref{ch:cmp:physical-scalar-bound} and~\eqref{ch:cmp:field-square}
therefore imply
\[
 R_{\widetilde g}\ge2e^{-4s_0}|\alpha|_{\bar g}^2
       \ge2e^{-8s_0}|\alpha|_{\bar g}^2
       =2\sum_{i=1}^2|\widetilde X_i|_{\widetilde g}^2.
\]
The last equality uses the isometry
$X\mapsto\iota_XdV_{\widetilde g}$ between vectors and two-forms in
dimension three.  Closedness gives the divergence equations because
$d(\iota_XdV_{\widetilde g})=(\Div_{\widetilde g}X)dV_{\widetilde g}$.
Also $\widetilde g=e^{4s_0}(g+l^2df\otimes df)\ge g$.
A $\widetilde g$-Cauchy sequence is therefore $g$-Cauchy; its limit lies in
$X$ with $B$ included, and local smooth comparison of the two metrics
gives convergence for $\widetilde g$ as well.  Completeness follows, and
the boundary sign is Lemma~\ref{ch:cmp:boundary}.
\end{proof}

The asymptotic estimates, scalar-curvature integrability and ADM energy
control needed for the final Riemannian comparison are established in
Proposition~\ref{ch:mass:comparison-energy}.
